\section{从按出虎水到黄河：征服不是一条直线}

\subsection{一一一四年：东北的道路先被切断}

按出虎水流入阿什河，沿河的营地、寨堡和渡口并不属于一张空白地图。辽朝在东北经营女真诸部，既依赖军政首领，也依赖能够输送皮货、粮食和人力的地方关系。完颜阿骨打起兵时，所争的首先是这些关系：谁能召集骑从，谁能守住河谷通路，谁能使邻近部众相信辽的惩罚已经不能恢复旧秩序。后来史书把胜利写入帝王本纪，容易使读者以为一个国号从一开始便统摄了东北。实际上，起兵使原来可以回避的选择变得急迫，部众、城寨和辽方守军才被迫在不断变化的形势中表明去向。

阿骨打一一一五年称帝建金。这个日期可以由《金史·太祖纪》与辽方纪年相互核对；它能确认新朝廷公开建立，却不能量化所有部众的服从。称帝把军事首领、赏赐、征发和对外使节集中到一个可见中心，也把战争从局部反抗变成与辽争夺城镇、道路和人力的竞争。金的早期动员不能只归于骑兵技艺。骑乘者也要有马匹、草场、冬季驻地、修补器械的人和传递命令的线路；一支离开东北河谷的军队，越依赖既有城市和仓储，越必须学会同不同语言、不同职业的人群打交道。

《金史》叙述建国时多从君长和战绩着眼，《辽史》则将女真行动安放在辽末危机中。两书都保存了重要的次序，却都不擅长记录普通家庭怎样理解征兵、贸易中断或逃亡。金代碑志和东北城址可以把少数官名、地名与营建活动落到实处；它们反过来也提醒我们，能够留下碑刻者本已不是全部居民。因而，本章不把“女真”当成一个行动一致的主体。渤海后裔、契丹人、奚人、汉人以及诸部中不同的首领、士兵和家户，进入新秩序的时间和条件不可能相同。

\subsection{海上之盟的海并不通向同一个目标}

宋朝使者经登州渡海，是因为辽仍隔在宋与女真政权之间；海路缩短了外交距离，也隔开了双方对燕云的想象。宋廷希望借新盟友取得失去已久的边地，金方则要借宋的财货、情报和牵制集中对付辽。双方写下的承诺并未消除这种差别。盟约可以规定名义、岁币或攻守方向，却无法预先决定攻城之后谁来接收仓库、谁来安顿降人、谁来负担守城的粮草。

一一一八年起，宋金围绕攻辽交涉，到一一二〇年前后形成通常所称的海上之盟。宋方编年材料详记朝廷议论，金方史书更多强调本朝战事；二者并读只能说明交涉确曾把燕云推到共同议题上。条款的准确文字、使者所能代表的范围以及地方对未来统治者的态度，都不能由后出的叙事补齐。对住在边地的人而言，新的盟约很可能先表现为军队过境、粮秣征发和市场消息，而不是一段可以当场阅读的外交文字。

辽南京被金军攻下后，燕京成为试验盟约的地方。金把部分地区交给宋，同时要求以岁币和边防条件维系安排。宋军接到的不是未曾经营的土地，而是刚经历战斗、转运和人口扰动的城与州县。城门的归属固然醒目，城外的粮田、驿道、河道和守军给养却决定接收能否延续。宋方在燕地的困难让金方看见其动员的限度；金方对燕云的处置又使宋廷意识到，盟友并不会替它承担恢复旧边界的全部成本。辽的覆亡没有带来缓冲，反而让两个原先隔海相望的政权直接面对。

\subsection{从燕云到汴京：推进需要新的粮道}

一一二五年辽天祚帝被俘，辽朝覆亡。金军随即转向北宋。这一连串行动常被压缩为“金灭辽、金灭宋”，但两次战争的空间条件很不相同。燕云以北的行军仍可连接金的早期人力和军政网络；进入河北、黄河与开封方向，军队面对的是人口更密集、城市更多、漕运和仓储更复杂的地区。攻城的胜利只解决了一时的军事问题，驻军、搜集粮草、看守俘虏、传递诏令和防止后路断绝，则把战争拉入地方日常。

围汴京的材料最能显示叙事立场。宋方记录中的城陷、议和与掳掠带着亡国的痛感；《金史》则把南下放入本朝开拓的节奏。两者都不能照单全收，尤其不能以相互竞争的军数、斩获和个人动机填满现场。较稳妥的判断是：一一二六年的围城与一一二七年的城陷，靖康城陷与北宋覆亡，改变了北方最高政治中心的归属，也把大量官员、工匠、宗室和居民卷入迁徙、俘虏和重新编户的过程。至于每一种人经历了什么，保存下来的材料远比后来的道德判断零碎。

汴京的陷落不能证明所有华北立刻进入同一种统治。城内的宫廷、官署和仓储与乡村、盐场、渡口的变化速度不会一致；一个地方宣布归附，也不等于税粮、司法和征兵已经按新制度运行。金廷若想保有南下的成果，必须决定保留哪些旧官、如何处置旧税源、以何种方式让军队获得供给。征服在这里转成了行政问题。它也给南方留下一个长期压力：逃至江淮和东南的人群带去信息、技艺和政治合法性的争执，南宋的建立不只是一次皇位延续，也是在北方失序中重新组织道路、军费和人口的结果。

\subsection{刘豫并非一块可以绕开的插曲}

金军继续向江淮推进，发现直接占据广阔中原的代价极高。扶立刘豫政权，是把一部分征税、安抚、任官与防宋的风险交给一个可被利用的中介。这样的安排并不表示金完全放手，也不表示刘豫政权只是没有社会基础的木偶。它要求地方官、军队、乡里和商人每天在多重命令之间作出选择；有人借此保存职位和财产，有人离开，有人等待局势再变。材料主要记住了朝廷的废立，较少记录县乡的协商，因此不能把一一三一年后的中原生活写成整齐的服从或整齐的抵抗。

一一三七年金废刘豫，转行直接治理。这项决定显示代理统治不能稳定解决军费、任官和边防的难题。直接经营要求更多文书、官员与运输线，也要求金廷承认旧宋社会并不会因一次政令而重新分类。猛安谋克为女真军政组织提供了重要骨架，州县和路的行政传统又承担了大量地方事务；二者在不同地区如何交错，不能由单一制度名推断。恰在这种不均衡中，金从东北兴起的军事政权逐步成为经营华北的王朝。

\begin{debate}{征服者与被征服者之间没有一条固定界线}
后世常把金的统治写成女真与汉人两种完整社会的对峙。金史官修叙事、南宋对手文字和近现代分类都可能强化这种二分。碑志中的官衔、家族迁徙和多种文字材料显示的却是更具体的地域、职业和政治关系。猛安谋克、州县、寺院、市场与军镇所形成的身份并不相同；本章不以姓名、服饰或一通碑刻替任何人规定稳定的“族群立场”。
\end{debate}
